\textbf{Desarrollo de las magnitudes intermedias.} Con $D=AV^2+Bn^2/V^2$, el empuje disponible limita la carga a $n_T=\sqrt{(T-AV^2)V^2/B}$. La sustentación la limita a $n_L=qSC_{L,max}/W$. El valor estructural 3.2 es un dato del enunciado, no un cálculo; se elige el menor de los tres límites.\begin{enumerate}

\item Peso del estado.

\[\fitmath{W=6200\cdot 9.80665\approx 60801.23\,\mathrm{N}}\]

\item Fuerza de referencia local.

\[\fitmath{\mathrm{qS}=0.5\cdot 0.7361155474\cdot {110}^{2}\cdot 34\approx 151419\,\mathrm{N}}\]

\item Coeficiente parásito.

\[\fitmath{A=0.5\cdot 0.7361155474\cdot 34\cdot 0.023\approx 0.2878212\,\mathrm{kg/m}}\]

\item Coeficiente cuadrático de carga.

\[\fitmath{B=\frac{2\cdot 0.043\cdot {W}^{2}}{0.7361155474\cdot 34}}\]

\[\fitmath{B=\frac{2\cdot 0.043\cdot {(60801.23)}^{2}}{0.7361155474\cdot 34}\approx 1.270277\times10^{7}\,\mathrm{N\,m^{2}/s^{2}}}\]

\item Carga que agota el margen de empuje.

\[\fitmath{\mathrm{nT}=\sqrt{\frac{\left(8200-A\cdot {110}^{2}\right)\cdot {110}^{2}}{B}}}\]

\[\fitmath{\mathrm{nT}=\sqrt{\frac{\left(8200-(0.287821179)\cdot {110}^{2}\right)\cdot {110}^{2}}{(12702765.3)}}\approx 2.119792\,}\]

\item Carga que alcanza CL máximo.

\[\fitmath{\mathrm{nL}=\frac{\mathrm{qS}\cdot 1.6}{W}}\]

\[\fitmath{\mathrm{nL}=\frac{(151418.9681)\cdot 1.6}{(60801.23)}\approx 3.984629\,}\]

\end{enumerate}
\par Comparar $n_T\simeq2.120$, $n_L\simeq3.985$ y el dato $3.2$: el límite activo es el empuje y la carga sostenida máxima del modelo es aproximadamente 2.120.
